\documentclass{article}
\usepackage[utf8]{inputenc}
\usepackage{geometry}
\geometry{a4paper, margin=1in}

\title{Respuestas - Evaluación Semana 1 (A/AAAA)}
\author{Gabriel}
\date{\today}

\begin{document}

\maketitle

\section*{Soluciones al Cuestionario}

\begin{enumerate}
    \item \textbf{Diferencia entre registros A y AAAA:}
    La diferencia fundamental radica en la versión del Protocolo de Internet (IP) que utilizan:
    \begin{itemize}
        \item Un registro \textbf{A} (Address) mapea un nombre de dominio a una dirección \textbf{IPv4} (ej. \texttt{192.0.2.1}).
        \item Un registro \textbf{AAAA} (Quad A) mapea un nombre de dominio a una dirección \textbf{IPv6} (ej. \texttt{2001:0db8:85a3:0000:0000:8a2e:0370:7334}).
    \end{itemize}

    \item \textbf{Propósito de \texttt{dns.resolver.resolve()}:}
    Es la función principal de la biblioteca \texttt{dnspython} para realizar consultas DNS. Su propósito es enviar una solicitud a un servidor DNS para "resolver" un nombre de dominio, pidiendo un tipo de registro específico (como 'A', 'AAAA', 'MX', etc.) y devolver los resultados encontrados.

    \item \textbf{Múltiples registros 'A':}
    Que un dominio devuelva múltiples registros 'A' significa que ese nombre de dominio está asociado a varias direcciones IP. Esta es una técnica común llamada \textbf{DNS Round Robin}, utilizada para el \textbf{balanceo de carga}. Los clientes que resuelven el dominio reciben la lista de IPs y pueden conectarse a cualquiera de ellas, distribuyendo así el tráfico entre varios servidores. Esto mejora la disponibilidad y la capacidad de respuesta de sitios web con alto tráfico.

\end{enumerate}

\end{document}